
\section{Motivação}

\frame{
  \frametitle{Ruby}
  \pause
  Linguagens, como qualquer outro programa, possuem versões!

  \vspace{1em}
  \begin{itemize}
    \pause
    \item No caso, Ruby tem lançamentos constantes
          \pause
          \begin{itemize}
            \item Só em 2015, foram 12 versões novas entre minors a patches
          \end{itemize}
    \pause
    \item Além disso, existem diferentes implementações
          \pause
          \begin{itemize}
            \item JRuby
            \item Rubinious
            \item etc.
          \end{itemize}
    \pause
    \item Como conciliar isso com um ambiente de desenvolvimento estável?
  \end{itemize}
}

\begin{frame}
  \frametitle{RubyGems}
  \begin{columns}
    \column[c]{.4\textwidth}
      \begin{center}
        \includegraphics[width=0.8\textwidth]{images/favicon}
      \end{center}
    \column[c]{.6\textwidth}
      \pause
      ``Bibliotecas'' do Ruby

      \vspace{1em}
      \begin{itemize}
        \pause
        \item Cada gem tem uma versão
        \pause
        \item Gems dependem de outras gems
        \pause
        \item Diferentes para dev, testes e produção.
        \pause
        \item RubyGems.org
      \end{itemize}
  \end{columns}

  \pause
  \vspace{2em}
  Mais um nível de complexidade no gerenciamento de versões em projetos!
\end{frame}
